% The host stack, and the two paths through it.
\begin{tikzpicture}[font=\sffamily\small, x=1cm, y=1cm]

% ---------------- the host stack
\node[user, text width=34mm] (app) at (-5.6,4.0)
  {setpoint source\\{\scriptsize about twenty lines of standard library, no dependency}};
\node[user, text width=34mm] (sock) at (-5.6,1.9)
  {a raw socket\\{\scriptsize data frames and error frames arrive on the same one}};
\node[kern, text width=34mm] (net) at (-5.6,-0.2)
  {the kernel: a network device\\{\scriptsize link state, counters, filters, timestamps, automatic restart}};
\begin{scope}[on background layer]
\node[layer, fit=(app)(sock)(net), inner sep=6pt,
      label={[layerlabel, anchor=south west]north west:the motion master}] {};
\end{scope}

% ---------------- the two paths out
\node[kern, text width=30mm] (vcan) at (0.2,1.9)
  {the virtual interface\\{\scriptsize real frames, real kernel path, no hardware at all}};
\node[hw, text width=30mm] (drv) at (0.2,-0.2)
  {the driver and the adapter\\{\scriptsize one of three kinds, and only one of them works here}};

% ---------------- the three adapters
\node[absentblk, text width=32mm] (a1) at (5.6,1.9)
  {a hat with a classic controller\\{\scriptsize cheap, and it objects to flexible-data traffic rather than ignoring it}};
\node[absentblk, text width=32mm] (a2) at (5.6,-0.4)
  {a closed-library USB device\\{\scriptsize its own frame structure, no licence file, and the standard tools cannot see it}};
\node[hw, text width=32mm] (a3) at (5.6,-2.7)
  {an adapter presenting the kernel's interface\\{\scriptsize permissively licensed, genuine flexible-data rates. Chosen}};

% ---------------- the bus
\node[busnode, text width=26mm] (node) at (0.2,-3.4)
  {the joint node\\{\scriptsize chapter 9}};

\draw[flow] (app) -- (sock);
\draw[flow] (sock) -- (net);
\draw[bus] (net) -- node[lbl, above]{no hardware} (vcan);
\draw[bus] (net) -- (drv);
\draw[hiz, ->] (drv) -- (a1);
\draw[hiz, ->] (drv) -- (a2);
\draw[flow] (drv) -- (a3);
\draw[bus, line width=1.1pt] (a3.south) -- ++(0,-0.7) -| node[lbl, above, pos=0.25]{the wire} (node.east);

\node[chlabel, anchor=west] at (2.0,2.9) {chapter 20's tests run here};

\node[umlnote, text width=54mm, anchor=north west] at (-8.4,-2.2)
  {The upper path is the whole first half of the chapter: two commands create an
   interface that carries real frames through the real kernel code, so the
   socket, the filters, the tools, the recording format and the test harness are
   all exercised before any adapter is bought. What it cannot teach is timing,
   because frames on it have no bit timing and no arbitration at all.};
\end{tikzpicture}
